\documentclass[12pt]{article}
\usepackage[margin=1in]{geometry}
\usepackage[T1]{fontenc}
\usepackage{xcolor}
\definecolor{garnet}{HTML}{73000A}

\setlength{\parindent}{0pt}
\setlength{\parskip}{0.3em}

\begin{document}

\noindent
Department of Mechanical Engineering\\
University of South Carolina\\
Columbia, SC 29208\\
\today

\vspace{0.5em}

\noindent
SC Space Grant Consortium\\
College of Charleston\\
Charleston, SC 29424

\vspace{0.5em}

\noindent
\textbf{Re: Faculty Recommendation for Noah Young --- Undergraduate Research Award}

\noindent
{\color{garnet}\rule{\textwidth}{0.5pt}}

\vspace{0.3em}

\noindent
Dear Members of the Evaluation Committee:

I recommend Noah Young for the SC Space Grant Undergraduate Research Award. As his faculty advisor in the Integrated Multiphysics and Systems Engineering Laboratory (iMSEL) since August 2025, I have supervised his research on radar simulation and autonomous perception systems from his first week of college. I have supervised dozens of undergraduate researchers over my career. Among all undergraduates, Noah's output places him solidly above average; among freshmen, he is in the top quartile. Measured against the broader pool of students I have taught, he ranks in the top five percent.

Noah joined iMSEL at the start of his first semester as an Honors College aerospace engineering student. He is now contributing to an ASME IDETC 2026 submission on synthetic radar data augmentation, in which he developed a method using land/water segmentation and object extraction to regenerate radar imagery from real sensor data. Radar fundamentals---signal geometry, terrain interaction, image processing---were entirely new to him. Rather than waiting for instruction, he independently read background literature on maritime radar captures, FMCW radar, and geometric simplification methods, then applied what he learned to produce working code within weeks. This self-directed learning style is characteristic of how he operates: he would rather teach himself a topic first and arrive with informed questions than consume lab time on basics. He maintains a 4.0\,GPA while carrying this research workload.

His technical contributions span multiple systems in our lab, and what distinguishes him is the level of initiative and polish he brings to each task. He fine-tuned a YOLO model for maritime small-object detection and deployed it on the NVIDIA Jetson Orin. He designed a 3D-printable radar antenna mount in Creo Parametric that enabled our Furuno NXT radar to be integrated onto a research vessel for offshore data collection. He built the raw radar-to-PNG conversion pipeline---choosing Rust for its performance with thousands of files---and the real-time visualization tools that our entire team now uses for experimental work. In each case, he did not simply complete the assigned task; he identified adjacent problems (file processing speed, export compatibility, visualization format) and solved them without being asked. He is consistently the fastest member of the lab to respond to emails, and he visibly puts more effort into the quality of his deliverables than many undergraduates who treat research as a job rather than an opportunity.

Noah delivers weekly technical presentations and monthly written reports to our Department of Defense sponsors. His monthly reports are organized with explicit objective statements, challenge-and-resolution narratives, and structured future-work sections---a format he adopted on his own, not one we prescribed. For a freshman explaining radar signal processing to program managers who are not domain specialists, this level of structured communication is uncommon. Most undergraduates at this stage summarize what they did; Noah frames why it matters and what comes next.

His proposed project---an open-source terrain-occluded radar simulator validated against measured returns---grows directly from his current work building the radar data visualization and conversion tools. He has already demonstrated that he can take an unfamiliar domain (radar physics), learn it independently, and produce lab-infrastructure tools that the full team relies on. The proposed simulator extends this trajectory: he would move from visualizing real radar data to generating synthetic data and validating it against our existing offshore captures. The real data exists in our lab, and I have reviewed his validation framework. The project is scoped appropriately for an undergraduate researcher who has already proven he can deliver functional software on schedule.

Noah is an AIAA member, participates in the USC Fly Club, and holds a MathWorks certification in Simscape. He is proficient in Python, C++, Rust, and MATLAB.

I recommend Noah Young for this award.

\vspace{0.8em}

\noindent
Sincerely,

\vspace{1.5em}

\noindent
Dr.\ Yi Wang\\
Associate Professor\\
Department of Mechanical Engineering\\
University of South Carolina\\
\textit{ywang@sc.edu}

\end{document}
